% Ad Familiares 4.15 --- headnote
% ad-familiares-04-15, August 46 BC (shortly after Fam. 4.14), Rome

\textit{Cicero to Cn. Plancius, in exile at Corcyra, written at
Rome shortly after the preceding letter (4.14) --- Perseus:
\textit{Romae paulo post ep. xiv a. 708 (46)}. Plancius, the
former tribune and aedile who in 54 BC had been the subject of
Cicero's surviving speech \textit{Pro Plancio}, had remained
loyal to Pompeius in the civil war and was now living, with
many fellow Pompeians, in proscribed exile at Corcyra. He had
sent Cicero a very short letter, which carried his affection
intact but said almost nothing about how he was bearing his
condition.}

\textit{Cicero's reply is itself only one of the shortest pieces
in the corpus, and it is built on the small distinction between
the two things a friend's letter from exile can tell you ---
that he loves you, and that he is enduring. Plancius's note had
made the first plain and left the second a blank, and Cicero
will not let the blank stand. The note's centre of gravity is
the opposition between \textit{propria fortuna} and \textit{in
communi sumus}: there is no private misery here, only a shared
one, and the consolation lies precisely in refusing to claim a
peculiar share of it. The closing antithesis --- ``which on your
side I can hope for, and on mine I can guarantee'' --- is the
characteristic gesture of these short letters of 46, in which
Cicero offers his own steadiness as the warrant for asking
steadiness in return.}
